%================================================
%    sheaves on manifolds 演習ノート
%================================================
\title{Exercises of \textit{Sheaves on Manifolds}}
\author{Toshi2019}
\date{\today 更新}

% -----------------------
% preamble
% -----------------------
% ここから本文 (\begin{document}) までの
% ソースコードに変更を加えた場合は
% 編集者まで連絡してください. 
% Don't change preamble code yourself. 
% If you add something
% (usepackage, newtheorem, newcommand, renewcommand),
% please tell it 
% to the editor of institutional paper of RUMS.

% ------------------------
% documentclass
% ------------------------
\documentclass[11pt, a4paper, dvipdfmx, leqno]{jsreport}

% ------------------------
% usepackage
% ------------------------
\usepackage{algorithm}
\usepackage{algorithmic}
\usepackage{amscd}
\usepackage{amsfonts}
\usepackage{amsmath}
\usepackage[psamsfonts]{amssymb}
\usepackage{amsthm}
\usepackage{ascmac}
\usepackage{bm}
\usepackage{color}
\usepackage{enumerate}
\usepackage{fancybox}
\usepackage[stable]{footmisc}
\usepackage{graphicx}
\usepackage{listings}
\usepackage{mathrsfs}
\usepackage{mathtools}
\usepackage{otf}
\usepackage{pifont}
\usepackage{proof}
\usepackage{subfigure}
\usepackage{tikz}
\usepackage{verbatim}
\usepackage[all]{xy}
\usepackage{url}
\usetikzlibrary{cd}



% ================================
% パッケージを追加する場合のスペース 
%\usepackage{calligra}
\usepackage[dvipdfmx]{hyperref}
\usepackage{xcolor}
\definecolor{darkgreen}{rgb}{0,0.45,0} 
\definecolor{darkred}{rgb}{0.75,0,0}
\definecolor{darkblue}{rgb}{0,0,0.6} 
\hypersetup{
    colorlinks=true,
    citecolor=darkgreen,
    linkcolor=darkred,
    urlcolor=darkblue,
}
\usepackage{pxjahyper}
%\usepackage[mathcal]{euscript}
\usepackage{layout}
\usepackage{framed}
\definecolor{lightgray}{rgb}{0.75,0.75,0.75}
\renewenvironment{leftbar}{%
  \def\FrameCommand{\textcolor{lightgray}{\vrule width 0.7zw} \hspace{10pt}}% 
  \MakeFramed {\advance\hsize-\width \FrameRestore}}%
{\endMakeFramed}
\newenvironment{redleftbar}{%
  \def\FrameCommand{\textcolor{red}{\vrule width 1pt} \hspace{10pt}}% 
  \MakeFramed {\advance\hsize-\width \FrameRestore}}%
 {\endMakeFramed}

%=================================

%============================================================
% layout for a4 9pt
%\setlength{\textwidth}{459pt}
%\setlength{\marginparwidth}{105pt}
%\setlength{\evensidemargin}{69pt}
%============================================================

% --------------------------
% theoremstyle
% --------------------------
\theoremstyle{definition}

% --------------------------
% newtheoem
% --------------------------

% 日本語で定理, 命題, 証明などを番号付きで用いるためのコマンドです. 
% If you want to use theorem environment in Japanece, 
% you can use these code. 
% Attention!
% All theorem enivironment numbers depend on 
% only section numbers.
\newtheorem{Axiom}{公理}[section]
\newtheorem{Definition}[Axiom]{定義}
\newtheorem{Theorem}[Axiom]{定理}
\newtheorem{Proposition}[Axiom]{命題}
\newtheorem{Lemma}[Axiom]{補題}
\newtheorem{Corollary}[Axiom]{系}
\newtheorem{Example}[Axiom]{例}
\newtheorem{Claim}[Axiom]{主張}
\newtheorem{Property}[Axiom]{性質}
\newtheorem{Attention}[Axiom]{注意}
\newtheorem{Question}[Axiom]{問}
\newtheorem{Problem}[Axiom]{問題}
\newtheorem{Consideration}[Axiom]{考察}
\newtheorem{Alert}[Axiom]{警告}
\newtheorem{Fact}[Axiom]{事実}
\newtheorem{com}[Axiom]{コメント}


% 日本語で定理, 命題, 証明などを番号なしで用いるためのコマンドです. 
% If you want to use theorem environment with no number in Japanese, You can use these code.
\newtheorem*{Axiom*}{公理}
\newtheorem*{Definition*}{定義}
\newtheorem*{Theorem*}{定理}
\newtheorem*{Proposition*}{命題}
\newtheorem*{Lemma*}{補題}
\newtheorem*{Example*}{例}
\newtheorem*{Corollary*}{系}
\newtheorem*{Claim*}{主張}
\newtheorem*{Property*}{性質}
\newtheorem*{Attention*}{注意}
\newtheorem*{Question*}{問}
\newtheorem*{Problem*}{問題}
\newtheorem*{Consideration*}{考察}
\newtheorem*{Alert*}{警告}
\newtheorem*{Fact*}{事実}
\newtheorem*{com*}{コメント}



% 英語で定理, 命題, 証明などを番号付きで用いるためのコマンドです. 
%　If you want to use theorem environment in English, You can use these code.
%all theorem enivironment number depend on only section number.
\newtheorem{Axiom+}{Axiom}[section]
\newtheorem{Definition+}[Axiom+]{Definition}
\newtheorem{Theorem+}[Axiom+]{Theorem}
\newtheorem{Proposition+}[Axiom+]{Proposition}
\newtheorem{Lemma+}[Axiom+]{Lemma}
\newtheorem{Example+}[Axiom+]{Example}
\newtheorem{Corollary+}[Axiom+]{Corollary}
\newtheorem{Claim+}[Axiom+]{Claim}
\newtheorem{Property+}[Axiom+]{Property}
\newtheorem{Attention+}[Axiom+]{Attention}
\newtheorem{Question+}[Axiom+]{Question}
\newtheorem{Problem+}[Axiom+]{Problem}
\newtheorem{Consideration+}[Axiom+]{Consideration}
\newtheorem{Alert+}{Alert}
\newtheorem{Fact+}[Axiom+]{Fact}
\newtheorem{Remark+}[Axiom+]{Remark}

% ----------------------------
% commmand
% ----------------------------
% 執筆に便利なコマンド集です. 
% コマンドを追加する場合は下のスペースへ. 

% 集合の記号 (黒板文字)
\newcommand{\NN}{\mathbb{N}}
\newcommand{\ZZ}{\mathbb{Z}}
\newcommand{\QQ}{\mathbb{Q}}
\newcommand{\RR}{\mathbb{R}}
\newcommand{\CC}{\mathbb{C}}
\newcommand{\PP}{\mathbb{P}}
\newcommand{\KK}{\mathbb{K}}


% 集合の記号 (太文字)
\newcommand{\nn}{\mathbf{N}}
\newcommand{\zz}{\mathbf{Z}}
\newcommand{\qq}{\mathbf{Q}}
\newcommand{\rr}{\mathbf{R}}
\newcommand{\cc}{\mathbf{C}}
\newcommand{\pp}{\mathbf{P}}
\newcommand{\kk}{\mathbf{K}}

% 特殊な写像の記号
\newcommand{\ev}{\mathop{\mathrm{ev}}\nolimits} % 値写像
\newcommand{\pr}{\mathop{\mathrm{pr}}\nolimits} % 射影

% スクリプト体にするコマンド
%   例えば　{\mcal C} のように用いる
\newcommand{\mcal}{\mathcal}

% 花文字にするコマンド 
%   例えば　{\h C} のように用いる
\newcommand{\h}{\mathscr}

% ヒルベルト空間などの記号
\newcommand{\F}{\mcal{F}}
\newcommand{\X}{\mcal{X}}
\newcommand{\Y}{\mcal{Y}}
\newcommand{\Hil}{\mcal{H}}
\newcommand{\RKHS}{\Hil_{k}}
\newcommand{\Loss}{\mcal{L}_{D}}
\newcommand{\MLsp}{(\X, \Y, D, \Hil, \Loss)}

% 偏微分作用素の記号
\newcommand{\p}{\partial}

% 角カッコの記号 (内積は下にマクロがあります)
\newcommand{\lan}{\langle}
\newcommand{\ran}{\rangle}



% 圏の記号など
\newcommand{\Set}{{\bf Set}}
\newcommand{\Vect}{{\bf Vect}}
\newcommand{\FDVect}{{\bf FDVect}}
\newcommand{\Mod}{\mathop{\mathrm{Mod}}\nolimits}
\newcommand{\CGA}{{\bf CGA}}
\newcommand{\GVect}{{\bf GVect}}
\newcommand{\Lie}{{\bf Lie}}
\newcommand{\dLie}{{\bf Liec}}



% 射の集合など
\newcommand{\Map}{\mathop{\mathrm{Map}}\nolimits}
\newcommand{\Hom}{\mathop{\mathrm{Hom}}\nolimits}
\newcommand{\End}{\mathop{\mathrm{End}}\nolimits}
\newcommand{\Aut}{\mathop{\mathrm{Aut}}\nolimits}
\newcommand{\Mor}{\mathop{\mathrm{Mor}}\nolimits}

% その他便利なコマンド
\newcommand{\dip}{\displaystyle} % 本文中で数式モード
\newcommand{\e}{\varepsilon} % イプシロン
\newcommand{\dl}{\delta} % デルタ
\newcommand{\pphi}{\varphi} % ファイ
\newcommand{\ti}{\tilde} % チルダ
\newcommand{\pal}{\parallel} % 平行
\newcommand{\op}{{\rm op}} % 双対を取る記号
\newcommand{\lcm}{\mathop{\mathrm{lcm}}\nolimits} % 最小公倍数の記号
\newcommand{\Probsp}{(\Omega, \F, \P)} 
\newcommand{\argmax}{\mathop{\rm arg~max}\limits}
\newcommand{\argmin}{\mathop{\rm arg~min}\limits}





% ================================
% コマンドを追加する場合のスペース 

\makeatletter
\renewenvironment{proof}[1][\proofname]{\par
  \pushQED{\qed}%
  \normalfont \topsep6\p@\@plus6\p@\relax
  \trivlist
  \item[\hskip\labelsep
%        \itshape
         \bfseries
%    #1\@addpunct{.}]\ignorespaces
    {#1}]\ignorespaces
}{%
  \popQED\endtrivlist\@endpefalse
}
\makeatother

\renewcommand{\proofname}{証明.}


%\renewcommand\proofname{\bf 証明} % 証明
\numberwithin{equation}{section}
\newcommand{\cTop}{\textsf{Top}}
%\newcommand{\cOpen}{\textsf{Open}}
\newcommand{\Op}{\mathop{\textsf{Op}}\nolimits}
\newcommand{\Ob}{\mathop{\textrm{Ob}}\nolimits}
\newcommand{\id}{\mathop{\mathrm{id}}\nolimits}
\newcommand{\pt}{\mathop{\mathrm{pt}}\nolimits}
\newcommand{\Int}{\mathop{\mathrm{Int}}\nolimits}
\newcommand{\res}{\mathop{\rho}\nolimits}
\newcommand{\A}{\mcal{A}}
\newcommand{\B}{\mcal{B}}
\newcommand{\C}{\mcal{C}}
\newcommand{\D}{\mcal{D}}
\newcommand{\E}{\mcal{E}}
\newcommand{\G}{\mcal{G}}
%\newcommand{\H}{\mcal{H}}
\newcommand{\I}{\mcal{I}}
\newcommand{\J}{\mcal{J}}
\newcommand{\OO}{\mcal{O}}
\newcommand{\Ring}{\mathop{\textsf{Ring}}\nolimits}
\newcommand{\cAb}{\mathop{\textsf{Ab}}\nolimits}
\newcommand{\Ker}{\mathop{\mathrm{Ker}}\nolimits}
\newcommand{\im}{\mathop{\mathrm{Im}}\nolimits}
\newcommand{\rk}{\mathop{\mathrm{rank}}\nolimits}
\newcommand{\Coker}{\mathop{\mathrm{Coker}}\nolimits}
\newcommand{\Coim}{\mathop{\mathrm{Coim}}\nolimits}
\newcommand{\Ht}{\mathop{\mathrm{Ht}}\nolimits}
\newcommand{\supp}{\mathop{\mathrm{supp}}\nolimits}
\newcommand{\colim}{\mathop{\mathrm{colim}}}
\newcommand{\Ext}{\mathop{\mathrm{Ext}}\nolimits}
\newcommand{\Tor}{\mathop{\mathrm{Tor}}\nolimits}

\newcommand{\cat}{\mathscr{C}}

%筆記体
\newcommand{\cA}{\mcal{A}}
\newcommand{\cB}{\mcal{B}}
\newcommand{\cC}{\mcal{C}}
\newcommand{\cD}{\mcal{D}}
\newcommand{\cE}{\mcal{E}}
\newcommand{\cF}{\mcal{F}}
\newcommand{\cG}{\mcal{G}}
\newcommand{\cH}{\mcal{H}}
\newcommand{\cI}{\mcal{I}}
\newcommand{\cJ}{\mcal{J}}
\newcommand{\cK}{\mcal{K}}
\newcommand{\cL}{\mcal{L}}
\newcommand{\cM}{\mcal{M}}
\newcommand{\cN}{\mcal{N}}
\newcommand{\cO}{\mcal{O}}
\newcommand{\cP}{\mcal{P}}
\newcommand{\cQ}{\mcal{Q}}
\newcommand{\cR}{\mcal{R}}
\newcommand{\cS}{\mcal{S}}
\newcommand{\cT}{\mcal{T}}
\newcommand{\cU}{\mcal{U}}
\newcommand{\cV}{\mcal{V}}
\newcommand{\cW}{\mcal{W}}
\newcommand{\cX}{\mcal{X}}
\newcommand{\cY}{\mcal{Y}}
\newcommand{\cZ}{\mcal{Z}}

\newcommand{\rmD}{\mathrm{D}}
\newcommand{\vD}{\mathbf{D}}


\newcommand{\scA}{\mathscr{A}}
\newcommand{\scB}{\mathscr{B}}
\newcommand{\scC}{\mathscr{C}}
\newcommand{\scD}{\mathscr{D}}
\newcommand{\scE}{\mathscr{E}}
\newcommand{\scF}{\mathscr{F}}
\newcommand{\scN}{\mathscr{N}}
\newcommand{\scO}{\mathscr{O}}
\newcommand{\scR}{\mathscr{R}}
\newcommand{\scS}{\mathscr{S}}
\newcommand{\scV}{\mathscr{V}}

\newcommand{\ibA}{\mathop{\text{\textit{\textbf{A}}}}}
\newcommand{\ibB}{\mathop{\text{\textit{\textbf{B}}}}}
\newcommand{\ibC}{\mathop{\text{\textit{\textbf{C}}}}}
\newcommand{\ibD}{\mathop{\text{\textit{\textbf{D}}}}}
\newcommand{\ibE}{\mathop{\text{\textit{\textbf{E}}}}}
\newcommand{\ibF}{\mathop{\text{\textit{\textbf{F}}}}}
\newcommand{\ibG}{\mathop{\text{\textit{\textbf{G}}}}}
\newcommand{\ibH}{\mathop{\text{\textit{\textbf{H}}}}}
\newcommand{\ibI}{\mathop{\text{\textit{\textbf{I}}}}}
\newcommand{\ibJ}{\mathop{\text{\textit{\textbf{J}}}}}
\newcommand{\ibK}{\mathop{\text{\textit{\textbf{K}}}}}
\newcommand{\ibL}{\mathop{\text{\textit{\textbf{L}}}}}
\newcommand{\ibM}{\mathop{\text{\textit{\textbf{M}}}}}
\newcommand{\ibN}{\mathop{\text{\textit{\textbf{N}}}}}
\newcommand{\ibO}{\mathop{\text{\textit{\textbf{O}}}}}
\newcommand{\ibP}{\mathop{\text{\textit{\textbf{P}}}}}
\newcommand{\ibQ}{\mathop{\text{\textit{\textbf{Q}}}}}
\newcommand{\ibR}{\mathop{\text{\textit{\textbf{R}}}}}
\newcommand{\ibS}{\mathop{\text{\textit{\textbf{S}}}}}
\newcommand{\ibT}{\mathop{\text{\textit{\textbf{T}}}}}
\newcommand{\ibU}{\mathop{\text{\textit{\textbf{U}}}}}
\newcommand{\ibV}{\mathop{\text{\textit{\textbf{V}}}}}
\newcommand{\ibW}{\mathop{\text{\textit{\textbf{W}}}}}
\newcommand{\ibX}{\mathop{\text{\textit{\textbf{X}}}}}
\newcommand{\ibY}{\mathop{\text{\textit{\textbf{Y}}}}}
\newcommand{\ibZ}{\mathop{\text{\textit{\textbf{Z}}}}}

\newcommand{\ibx}{\mathop{\text{\textit{\textbf{x}}}}}

\newcommand{\Comp}{\mathop{\mathrm{C}}\nolimits}
\newcommand{\Komp}{\mathop{\mathsf{K}}\nolimits}
\newcommand{\Domp}{\mathsf{D}}
\newcommand{\Kompl}{\mathop{\mathsf{K}^\mathrm{+}}\nolimits}
\newcommand{\Kompu}{\mathop{\mathsf{K}^\mathrm{-}}\nolimits}
\newcommand{\Kompb}{\mathop{\mathsf{K}^\mathrm{b}}\nolimits}
\newcommand{\Dompl}{\mathop{\mathsf{D}^\mathrm{+}}\nolimits}
\newcommand{\Dompu}{\mathop{\mathsf{D}^\mathrm{-}}\nolimits}
\newcommand{\Dompb}{\mathop{\mathsf{D}^\mathrm{b}}\nolimits}
\newcommand{\Dbcon}{\mathop{\mathsf{D}^\mathrm{b}_{\rr^+}}\nolimits}
\newcommand{\Dbconic}{\mathop{\mathsf{D}^\mathrm{b}_{\rrp}}\nolimits}

\newcommand{\CCat}{\Comp(\cat)}
\newcommand{\KCat}{\Komp(\cat)}
\newcommand{\DCat}{\Domp(\cat)}%圏Cの複体のホモトピー圏
%\newcommand{\HOM}{\mathop{\mathscr{H}\hspace{-2pt}om}\nolimits}%内部Hom
\newcommand{\HOM}{\mathop{\mathcal{H}\hspace{-0.5pt}om}\nolimits}%内部Hom
\newcommand{\RHOM}{\mathop{\mathrm{R}\hspace{-1.5pt}\HOM}\nolimits}

\newcommand{\muS}{\mathop{\mathrm{SS}}\nolimits}
\newcommand{\RG}{\mathop{\mathrm{R}\hspace{-0pt}\Gamma}\nolimits}
\newcommand{\RHom}{\mathop{\mathrm{R}\hspace{-1.5pt}\Hom}\nolimits}
\newcommand{\Rder}{\mathrm{R}}

\newcommand{\simar}{\mathrel{\overset{\sim}{\longrightarrow}}}%内部Hom
\newcommand{\simra}{\mathrel{\overset{\sim}{\longleftarrow}}}%内部Hom

\newcommand{\hocolim}{{\mathrm{hocolim}}}
\newcommand{\indlim}[1][]{\mathop{\varinjlim}\limits_{#1}}
\newcommand{\sindlim}[1][]{\smash{\mathop{\varinjlim}\limits_{#1}}\,}
\newcommand{\Pro}{\mathrm{Pro}}
\newcommand{\Ind}{\mathrm{Ind}}
\newcommand{\prolim}[1][]{\mathop{\varprojlim}\limits_{#1}}
\newcommand{\sprolim}[1][]{\smash{\mathop{\varprojlim}\limits_{#1}}\,}

\newcommand{\Sh}{\mathrm{Sh}}
\newcommand{\PSh}{\mathrm{PSh}}

\newcommand{\Lloc}[1][]{\mathord{\mathcal{L}^1_{\mathrm{loc},{#1}}}}
\newcommand{\ori}{\mathord{\mathrm{or}}}
\newcommand{\Db}{\mathord{\mathscr{D}b}}
\newcommand{\Conti}{\mathord{\mathscr{C}}}

\newcommand{\mres}[2][]{{\left.{#1}\right\rvert}_{#2}}
\newcommand{\blk}{\mathord{\ \cdot\ }}

\newcommand{\tens}[1][]{\mathbin{\otimes_{\raise1.5ex\hbox to-.1em{}{#1}}}}
\newcommand{\ttens}[1][]{\mathbin{\mathop{\overset{\mathrm{}}{\tens}}_{#1}}}

\newcommand{\etens}{\mathbin{\boxtimes}}
%\newcommand{\ltens}[1][]{\mathbin{\overset{\mathrm{L}}\tens}_{#1}}
\newcommand{\ltens}[1][]{\mathbin{\tens}^\mathrm{L}_{#1}}
\newcommand{\mtens}[1][]{\mathbin{\overset{\mathrm{\mu}}\tens}_{#1}}
\newcommand{\lltens}[1][]{\mathbin{\mathop{\tens}\limits^{\mathrm{L}}_{#1}}}
%\newcommand{\letens}{\overset{\mathrm{L}}{\etens}}
\newcommand{\detens}{\underline{\etens}}
\newcommand{\ldetens}{\overset{\mathrm{L}}{\underline{\etens}}}
\newcommand{\dtens}[1][]{{\overset{\mathrm{L}}{\underline{\otimes}}}_{#1}}

\newcommand{\letens}[2][]{\mathbin{\mathop{\overset{\mathrm{L}}{\etens}_{#1}}_{#2\ \ }}}
\newcommand{\lletens}[1][]{\mathbin{\mathop{\overset{\mathrm{L}}{\etens}_{#1}}}}

\newcommand{\wgld}{\mathop{\mathrm{wgld}}\nolimits}
\newcommand{\gld}{\mathop{\mathrm{gld}}\nolimits}
\newcommand{\hd}{\mathop{\mathrm{hd}}\nolimits}

\newcommand{\mtan}[1][]{T\hspace{-1.75pt}{#1}}
\newcommand{\mtp}[1][]{{}^{t\!}{#1}} % transpose of the morphism
\newcommand{\mpb}[1][]{{#1}_{\pi}} % pull-back of the morphism
%\newcommand{\mpb}[1][]{{}^{t\!}{#1}'} % pull-back of the morphism
\newcommand{\mdiff}[1][]{{#1}_{d}} % pull-back of the morphism



\newcommand{\pb}[1][]{\mathbin{\mathop{\times}\limits_{#1}}}

\newcommand{\cotb}[1][]{T^{\ast}{#1}}
\newcommand{\conm}[2][]{T_{#1}^{\hspace{0.7pt}\ast}{#2}}
\newcommand{\norb}[2][]{T_{#1}{#2}}

\newcommand{\cotz}[1][]{T^{\ast}_{#1}{#1}}


% =================================



%================================================
% 自前の定理環境
%   https://mathlandscape.com/latex-amsthm/
% を参考にした
\newtheoremstyle{mystyle}%   % スタイル名
    {5pt}%                   % 上部スペース
    {5pt}%                   % 下部スペース
    {}%              % 本文フォント
    {}%                  % 1行目のインデント量
    {\bfseries}%                      % 見出しフォント
    {.}%                     % 見出し後の句読点
    {12pt}%                     % 見出し後のスペース
    {\thmname{#1}\thmnumber{ #2}\thmnote{{\hspace{2pt}\normalfont (#3)}}}% % 見出しの書式

\theoremstyle{mystyle}
\newtheorem{AXM}{公理}[section]
\newtheorem{DFN}[Axiom]{定義}
\newtheorem{THM}[Axiom]{定理}
\newtheorem*{THM*}{定理}
\newtheorem{STT}{主張}[chapter]
\newtheorem*{STT*}{主張}
\newtheorem{PRP}[Axiom]{命題}
\newtheorem{LMM}[Axiom]{補題}
\newtheorem*{LMM*}{補題}
\newtheorem{CRL}[Axiom]{系}
\newtheorem{EG}[Axiom]{例}
\newtheorem{CNV}{規約}[chapter]
\newtheorem{NTT}[Axiom]{記法}
\newtheorem{RMK}[Axiom]{注意}
\newtheorem{CMT}{コメント}[chapter]
\newtheorem*{CMT*}{コメント}
\newtheorem{EXC}{演習}[chapter]


% 定理環境ここまで
%====================================================

% ---------------------------
% new definition macro
% ---------------------------
% 便利なマクロ集です

% 内積のマクロ
%   例えば \inner<\pphi | \psi> のように用いる
\def\inner<#1>{\langle #1 \rangle}

% ================================
% マクロを追加する場合のスペース 

%=================================





% ----------------------------
% documenet 
% ----------------------------
% 以下, 本文の執筆スペースです. 
% Your main code must be written between 
% begin document and end document.
% ---------------------------

\begin{document}
\maketitle
%\frontmatter
\layout
\chapter*{はじめに}
2025年度から始めた\cite{KS90}の演習セミナーのノート．

\section*{記号}
次の記号は断りなく使う．
\begin{itemize}
    \item 添字：
    なんらかの族$(a_i)_{i\in I}$を$(a_i)_i$とか$(a_i)$と
    略記することがある．
    \item 近傍：位相空間\(X\)の点\(x\)や部分集合\(Z\)に対し，
    その開近傍系をそれぞれ\(I_x\)や\(I_Z\)で表す．
    これらは，包含関係の逆で有向順序集合をなす．（\cite[Def1.1.7]{KS24b}の記法）
\end{itemize}
\tableofcontents
%\mainmatter
\chapter{ホモロジー代数}
%\setcounter{section}{2}
\section*{演習}
\setcounter{EXC}{16}
\begin{leftbar}
\begin{EXC}[ホモロジー次元]\label{EXC116}
    \(\cat\)をアーベル圏とする．
    \(n\)を非負自然数として，
    任意の\(X,Y\in\cat\)と\(j>n\)に対して\(\Ext^{j}(X,Y)=0\)であるとき，
    \(\cat\)はホモロジー次元は\(\leqq n\)であるという．
    ここで\(\Ext^{j}(X,Y)=\Hom_{\Domp(\cat)}(X,Y[j])\)である．
    \(\cat\)が十分多くの入射的対象をもつとする．
    次の条件は同値であることを示せ．    
    \begin{enumerate}[(i)]
        \item \(\cat\)のホモロジー次元は\(\leqq n\)である．
        \item 任意の\(X\in\cat\)に対し，
        \(X\)の長さ\(\leqq n\)の入射分解，
        すなわち完全列\(
            0\to X\to I^{0}\to \cdots\to I^{n}\to 0
        \)で各\(I^{j}\)が入射的であるものが存在する．
    \end{enumerate}
    
    これらの条件をみたす最小の\(
        n\in\nn\cup\{+\infty\}
    \)を\(\cat\)の\textbf{ホモロジー次元} (homological 
    dimension) とよび\(\hd(\cat)\)で表す．
\end{EXC}
\end{leftbar}


\setcounter{EXC}{27}

\begin{leftbar}
\begin{EXC}[大域ホモロジー次元]\label{EXC128}
    \(A\)を環とする．
    次の条件 (i)--(iii) が同値であることを示せ．    
    \begin{enumerate}[(i)]
        \item \(\Mod(A)\)のホモロジー次元は\(\leqq n\)である．
        \item 任意の左加群\(M\)が長さ\(\leqq n\)の入射分解をもつ．
        \item 任意の左加群\(M\)が長さ\(\leqq n\)の射影分解をもつ．\\（すなわち，完全列\(0\to P_{n}\to \cdots\to P_{0}\to M\to 0\)で各\(P_j\)が射影的であるものが存在する．）
    \end{enumerate}
    
    \(
        \gld(A)\coloneqq
        \sup\left(\hd(\Mod(A)),\hd(\Mod(A^{\op}))\right)
    \)とおき，\(\gld(A)\)を
    \(A\)の\textbf{大域ホモロジー次元} (global 
    homological dimension) とよぶ．
\end{EXC}
\end{leftbar}


\begin{leftbar}
\begin{EXC}[弱大域次元]\label{EXC129}
    \(A\)を環とする．
    \begin{enumerate}[(i)]
        \item 自由加群は射影加群であることを示せ．
        \item 射影加群は自由加群の自由加群の直和因子であることを示せ．
        \item 射影加群は平坦加群であることを示せ．
        \item \(n\geqq0\)を整数とする．
            次の条件(a)--\((\text{b})^{\op}\)が同値であることを示せ．
            \begin{enumerate}[(a)]
                \item 任意の\(j>n\), \(N\in\Mod(A^\op)\), 
                \(M\in\Mod(A)\)に対し，\(\Tor_j^A(N,M)=0\)
                \item 任意の\(M\in\Mod(A)\)に対し，分解\[
                    0\to P^n\to\dots\to P^0\to M\to0\quad\text{(\(P^j\)は平坦)}
                    \]が存在する．
                \item[\((\text{b})^{\op}\)]\setlength{\leftskip}{10pt}
                任意の\(M\in\Mod(A^\op)\)に対し，分解\[
                    0\to P^n\to\dots\to P^0\to M\to0\quad\text{(\(P^j\)は平坦)}
                \]が存在する．
            \end{enumerate}
            \(A\)の\textbf{弱大域次元} (weak global dimension) \(\wgld(A)\)を
            これらの条件を満たす最小の\(n\in\nn\cup\{+\infty\}\)として定める．
        \item \(\wgld(A)\leqq\gld(A)\)を示せ．
    \end{enumerate}
\end{EXC}
\end{leftbar}
\begin{proof}
    (i) 
    \(M\)を自由加群とする．左\(A\)加群の全射
    \(g\colon N\twoheadrightarrow N'\)に対し，\[
        g_\ast\colon\Hom_A(M,N)\to\Hom_A(M,N') 
        \quad\text{in \(\Mod(\zz)\)}
    \]が全射であることを示す．
    \(\psi\colon M\to N'\)を\(A\)加群の射とする．
    \(I\)を\(M\cong A^{\oplus I}\)となる添字集合とすると
    任意の\(m\in M\)は，
    \(M\)の生成系\((m_i)\)と\((a_i)_i\in A^{\oplus I}\)を
    用いて，\(m=\sum_{i\in I}a_im_i\)とかける．
    このとき，\[
        \psi(m)=\sum_{i}a_i\psi(m_i)\in N'
    \]であり，\(g\)が全射なので，\(n\in N\)で\[
        g(n)=\psi(m)=\sum_{i}a_i\psi(m_i),
        \quad
        \psi(m_i)=g(n_i)
    \]となるものがある．
    この\((n_i)_i\)に対して，\(\phi\colon M\to N\)を
    \[
        \phi(m_i)=n_i
    \]で定めると，
    \[
        \left(g_\ast(\phi)\right)(m_i)=g\circ\phi(m_i)=g(n_i)=\psi(m_i)
    \]となる．
    
    (ii) 
    \(P\)を射影加群とする．
    自由加群\(A^{\oplus I}\)と
    全射\(p\colon A^{\oplus I}\twoheadrightarrow P\)が存在する．
    実際，\(I= P\)として，
    \(p\)を\(p((a_x)_{x\in P})=\sum_{x\in P}a_xx\)と
    定めればよい．\(Q=\Ker p\)とすると，\[
        0\to Q\hookrightarrow A^{\oplus I}\twoheadrightarrow P\to0
    \]は完全列である．
    このとき，\(P\)が射影加群であることから，
    \(\id_P\)に対して，\(u\colon P\to A^{\oplus I}\)で
    \[p_\ast(u)=p\circ u=\id_P\]となる者が存在する．
    したがって，上の完全列は分裂し，
    \(A^{\oplus I}\cong P\oplus Q\)となる．
    
    (iii) 
    まず「自由\(\Rightarrow\)平坦」を示す．\(F=A^{\bigoplus I}\)を自由加群とし，\[
        0\to N_1\to N_2\to N_3\to0
    \]を右\(A\)加群の完全系列とする．\[
        0\to N_1\otimes A^{\bigoplus I}\to N_2\otimes A^{\bigoplus I}\to N_3\otimes A^{\bigoplus I}\to0
    \]において，\[
        N_1\otimes A^{\bigoplus I}\cong N_1^{\bigoplus I},\quad
        N_2\otimes A^{\bigoplus I}\cong N_2^{\bigoplus I}
    \]であり，\(j\colon N_1\to N_2\)は単射なので，\[
        \bigoplus_{i\in I}j_i\colon N_1^{\bigoplus I}\to N_2^{\bigoplus I}
    \]で
\end{proof}


%\clearpage
\chapter{層}\label{chap2}
\section*{演習}
\chapter{Poincar\'e-Verdier 双対性}\label{chap3}
\section*{演習}

\setcounter{EXC}{2}
\begin{leftbar}
\begin{EXC}\label{EXC3.3}
    \(X\)を位相空間とし，
    \(F\)を\(X\)上の層で\(\zz_X\)と局所的に同型であるものとする．
    次の標準同型が存在することを示せ．
    \[
        F\tens[]F\cong \zz_X,
        \quad
        \vD'{F}\cong F
    \]
\end{EXC}
\end{leftbar}

このノートだけの用語で演習\ref{EXC3.3}の性質を
\(F\)の対合律 (involution law) とよぶことにする．

まず「局所的に同型」の意味をはっきりさせる．
\begin{PRP}\label{prp:loc-isom-shv}
    \(F\)と\(G\)を\(X\)上の層とする．
    次の条件は同値である．
    \begin{enumerate}[(i)]
        \item 各点\(x\)に対し，開近傍\(U\)で\(\mres[F]{U}\cong\mres[G]{U}\)となるものが存在する．
        \item \(X\)の開被覆\((U_i)_{i\in I}\)で，各\(i\)に対し\(\mres[F]{U_{i}}\cong\mres[G]{U_{i}}\)となるものが存在する．
    \end{enumerate}
\end{PRP}
\begin{proof}
    (i)\(\Rightarrow\)(ii)：
    \(X\)の各点\(x\)に対し開近傍\(U_x\)で\(
        \mres[F]{U_{x}}\cong\mres[G]{U_{x}}
    \)となるものが存在する．
    \((U_x)_{x\in X}\)は\(X\)の開被覆である．

    (ii)\(\Rightarrow\)(i)：
    \((U_i)_{i\in I}\)を\(X\)の開被覆で
    各\(U_i\)に対し\(\mres[F]{U_{i}}\cong\mres[G]{U_{i}}\)となるものとする．
    \(x\in X\)とすると\(x\in U_{i}\)となる\(i\in I\)が存在する．
\end{proof}
\begin{DFN}
    \(X\)を位相空間とする．
    \(X\)上の層\(F\)と\(G\)が，
    命題\ref{prp:loc-isom-shv}の同値な条件をみたすとき，
    \(F\)と\(G\)は局所的に同型であるという．
\end{DFN}

演習\ref{EXC3.3}の証明にあたり，次の事実に注意しよう．
\begin{LMM}\label{lmm:pm1}
    \(\varphi\colon\zz\to\zz\)をアーベル群の同型とする．
    このとき\(\varphi\)は\(\pm\id_{\zz}\)のどちらか一方である．
\end{LMM}

\begin{proof}
    \(\varphi(1)=m\)とおくと
    任意の整数\(n\)に対し\(\varphi(n)=n\varphi(1)=nm\)となるので，
    \(\zz\)から\(\zz\)への射は\(m\)倍写像である．
    \(\psi\colon\zz\to\zz\)を\(\varphi\)の逆射とすると
    \[
        1=\psi(\varphi(1))=\psi(m)=m\psi(1)
    \]
    となる．\(\zz\)の可逆元は\(\pm1\)のみであるから，
    \[
        m=\pm1,
        \quad
        \psi(1)=\pm1
        \quad
        \text{（複合同順）}
    \]
    である．
\end{proof}

\begin{proof}[{演習\ref{EXC3.3}の証明．}]
    1つ目の主張を示す．
    \(X\)の開被覆\((U_i)_{i\in I}\)と各\(U_i\)上の層の射
    \[
        \theta_i\colon\mres[F]{U_{i}}\simar\mres[\zz_X]{U_{i}}
    \]で同型であるものが存在する．
    補題\ref{lmm:pm1}より，この\((\theta_i)_{i}\)は各\(i\)に対し
    \(\theta_i=\pm1\)とかける．実際，
    \(x\in U_i\)とすると，\(s_x\in F_x\cong\zz\)に対し\(
        (\theta_{i})_x(s_x)=\pm{s_x}
    \)となる．
    したがって，各\(i,j\)に対して    
    \begin{align*}
        \mres[\theta_{i}\otimes\theta_{i}]{{U_{i}\cap U_{j}}}
        &\cong
        (\pm1)\ttens[](\pm1)\\
        &\cong
        1\\
        &\cong
        \mres[\theta_{j}\otimes\theta_{j}]{{U_{i}\cap U_{j}}}
    \end{align*}
    となるので\((\theta_i\ttens[]\theta_i)_i\)から層の射
    \[
        \theta\tens[]\theta\colon F\tens[]F
        \to\zz_X\tens[]\zz_X
    \]
    がひきおこされる．これと同型
    \[
        \zz_X\tens[]\zz_X\to\zz_X
    \]
    の合成を\(\varphi\)で表す．
    各点\(x\in X\)に対して
    \[
        \varphi_x\colon F_x\ttens[]F_x\simar\zz\ttens[]\zz\cong(\zz_{X})_x
    \]が成り立つので，\(\varphi\)は同型である．

    1つ目の主張から2つ目の主張が従うことを示す．
    随伴\((\ttens[],\HOM)\)から
    \begin{align*}
        \Hom_{\zz_X}(F\ttens[]F,\zz_X)
        \cong
        \Hom_{\zz_X}(F,\HOM_{\zz_X}(F,\zz_X))
    \end{align*}
    である．
    同型は同型に送られるので同型\(
        F\ttens[]F\to \zz_X
    \)は同型\(F\to\HOM_{\zz_X}(F,\zz_X)=\rmD'F\)に送られる．
\end{proof}



\chapter{特殊化と超局所化}\label{chap4}
多様体\(X\)における，部分多様体\(M\)の法変形という幾何学的な構成法を復習したのち，
特殊化関手\(\nu_M\)という，\(X\)上の層に法束\(T_MX\)上の層を
対応させる関手と，超局所化関手\(\mu_M\)をそのフーリエ佐藤変換として定義する．
\(\mu_M\)の自然な一般化として\(\mu hom\)を定義し，
これらの関手の関手的性質を調べる．
\setcounter{CNV}{-1}
\begin{CNV}\label{cnv40}
    規約3.0はそのまま用いる．
    その上で，多様体とその間の射はすべて\(C^{\infty}\)級または実解析的であると仮定する．

    以降では，基本的に層の有界複体の導来圏で考える（cf.系3.12）．
    もちろん，結果の多くはもう少し弱い仮定の下でも成立する．
\end{CNV}
\section*{演習}
\chapter{層の超局所台}\label{chap5}
\section*{演習}
\chapter{超局所台と超局所化}\label{chap6}
\section*{演習}
\chapter{接触変換と純層}\label{chap7}
\section*{演習}
\chapter{構成可能層}\label{chap8}
\section*{演習}
\chapter{特性サイクル}\label{chap9}
\section*{演習}
\chapter{偏屈層}\label{chap10}
\section*{演習}
\chapter{\(\cO_X\)加群と\(\cD_X\)加群への応用}\label{chap11}
\section*{演習}
\appendix
\chapter{シンプレクティック幾何}\label{chapA}
\section*{演習}
























%\backmatter
%\appendix
%\chapter{aa}
%===============================================
% 参考文献スペース
%===============================================
\begin{thebibliography}{20} 
    \bibitem[B+84]{B+84} Borel, 
        \textit{Intersection Cohomology}, 
        Progress in Mathematics, 50, Birkh\"auser, 1984.
    \bibitem[Bou1]{Bou1} ブルバキ, 位相1, 東京図書, 1968.
    \bibitem[EGA IV]{EGA IV} Grothendieck, EGA IV, Pub. IHES., 32 1967.
    \bibitem[Fuk99]{Fuk99} 深谷賢治, シンプレクティック幾何学, 岩波講座 現代数学の展開 岩波書店, 1999.
    \bibitem[G58]{G58} Grauert, 
        \textit{On Levi's problem and the embedding of real analytic manifolds}, 
        Ann. Math. 68, 460--472 (1958).
    \bibitem[GP74]{GP74} Victor Guillemin, Alan Pollack, 
        \textit{Differential Topology}, 
        Prentice-Hall, 1974.
    \bibitem[KS85]{KS85} Masaki Kashiwara, Pierre Schapira, 
        \textit{Microlocal Study of Sheaves}, 
        Ast\`erisque 128, 1985.
    \bibitem[KS90]{KS90} Masaki Kashiwara, Pierre Schapira, 
        \textit{Sheaves on Manifolds}, 
        Grundlehren der Mathematischen Wissenschaften, 292, Springer, 1990.
    \bibitem[KS06]{KS06} Masaki Kashiwara, Pierre Schapira, 
        \textit{Categories and Sheaves}, 
        Grundlehren der Mathematischen Wissenschaften, 332, Springer, 2006.
    \bibitem[R55]{R55} de Rham, 
        \textit{Vari\'et\'es diff\'erentiables}, 
        Hermann, Paris, 1955.
    \bibitem[Sa59]{Sa59} Mikio Sato, 
        \textit{Theory of Hyperfunctions}, 
        1959--60.
    \bibitem[KS24b]{KS24b} Masaki Kashiwara, Pierre Schapira, 
        \textit{Introduction to Sheaves on Grothendieck Topologies}, 
        2024.
    \bibitem[S66]{S66} Schwartz, 
        \textit{Th\'eorie de distributions}, 
        Hermann, Paris, 1966.
    \bibitem[Sh16]{Sh16} 志甫淳, 層とホモロジー代数, 共立出版, 2016.
    \bibitem[Ike21]{Ike21} 池祐一, 超局所層理論概説, 2021.
    \bibitem[Tak17]{Tak17} 竹内潔, \(\D\)加群, 共立出版, 2017.
    %\bibitem[Og02]{Og02} 小木曽啓示, 代数曲線論, 朝倉書店, 2022.
\end{thebibliography}

%===============================================


\end{document}
